% FUNDAMENTAÇÃO TEÓRICA--------------------------------------------------------

\chapter{REFERENCIAL TEÓRICO}
\label{chap:referencial-teorico}

Demonstrar os assuntos relevantes à temática da pesquisa proposta, com base em publicações técnicas e científicas, sobretudo artigos científicos nacionais e internacionais, preferencialmente aqueles publicados nos últimos 10 anos. Deve contemplar, de maneira clara e objetiva, os elementos conceituais e a fundamentação teórica e científica que sejam importantes para respaldar a pesquisa, bem como servir de apoio para os procedimentos metodológicos inerentes ao tema proposto e para a interpretação dos resultados. 

As citações diretas longas (com mais de três linhas) devem constituir um parágrafo independente, recuado 4 cm da margem esquerda, justificado, fonte Arial ou Times New Roman, tamanho 10 e espaçamento 1 (simples) entre linhas, dispensando as aspas. É preciso indicar, obrigatoriamente, a fonte consultada, com o(s) nome(s) do(s) autor(es), ano de publicação e página. Caso o original contenha uma citação com uso de aspas, empregar aspas simples.

\begin{adjustwidth}{4cm}{}
    \SingleSpacing
    \small
    A NBR 10520 (ABNT, 2002b) recomenda, para as citações longas, o uso de recuo sem aspas; no entanto, quando houver necessidade de, no meio de uma citação longa, fazer-se uma interrupção para introduzir um comentário do autor, é preferível fechar a citação com aspas, fazer o comentário e abrir uma nova citação com aspas. (FRANÇA; VASCONCELLOS, 2013, p. 137).
    \vspace{0.4cm}
\end{adjustwidth}

\noindent \textcolor{cyan}{(No corpo do trabalho, tanto nas ilustrações como nas tabelas, a identificação aparece na parte superior. Na parte inferior, indica-se a fonte consultada – obrigatório, mesmo que seja produção de própria autoria conforme exemplo acima.)
\\
(A numeração das ilustrações deve ser sequencial (Figura \ref{fig:FIGURA1}, Figura \ref{fig:FIGURA2}), de acordo com a ordem que aparece no texto. Ao longo do texto, os títulos das ilustrações devem ser escritos da mesma forma e com os mesmos destaques gráficos (negrito, itálico, sublinhado, etc) que na lista de ilustrações. O mesmo acontece para as tabelas.)
}

\begin{figure}[!ht]
    \SingleSpacing
	\centering
	\caption{Nuvem de Tags dos PPCs das licenciaturas em Ciências Biológicas, Física e Química da UFRPE}
	\includegraphics[scale=1.5]{Figuras/Figura1.png}
    	\caption*{Fonte: Silva, 2019.}
	\label{fig:FIGURA1}
\end{figure}

Este capítulo também pode conter subdivisões conforme as peculiaridades de cada pesquisa, bem como a necessidade do autor. Neste caso, os subtítulos, também chamados de seções secundárias, devem ter numeração progressiva e seguindo o título principal; separado por um espaço 1,5 do texto que o precede e do que o sucede. Todos os itens precisam ser indicados no sumário.

\section{Título}
\label{sec:Título}

Texto texto texto texto texto texto texto texto texto texto texto texto texto texto texto texto texto texto texto texto texto texto texto texto texto texto texto texto texto.

Texto texto texto texto texto texto texto texto texto texto texto texto texto texto texto texto texto texto texto texto texto texto texto texto texto texto texto texto texto.

\subsection{Título}
\label{subsec:Título}

Texto texto texto texto texto texto texto texto texto texto texto texto texto texto texto texto texto texto texto texto texto texto texto texto texto texto texto texto texto texto texto.

\subsubsection{Título}
\label{subsubsec:Título}

Texto texto texto texto texto texto texto texto texto texto texto texto texto texto.

\begin{figure}[!ht]
    \SingleSpacing
	\centering
	\caption{Título título título}
	\includegraphics[scale=0.3]{Figuras/Figura2.PNG}
    	\caption*{Fonte: O autor, ano.}
	\label{fig:FIGURA2}
\end{figure}

Texto texto texto texto texto texto texto texto texto texto texto texto texto texto texto texto texto.

\subsubsubsection{Título}
\label{subsubsubsec:Título}


Texto texto texto texto texto texto texto texto texto texto texto texto texto texto texto texto texto.